\section{DAG-EVM}

\label{sec:dagevm}
%% ============================================================================

\subsection{Lineage: Block-STM Lifted}

Block-STM~\cite{blockstm} (Aptos Labs, 2022) is an optimistic concurrency
control (OCC) algorithm that speculatively executes a linear block's
transactions in parallel on separate copies of state, tracks per-transaction
read and write sets, and re-executes any transaction whose reads were
invalidated by an earlier-committed transaction's writes.

Block-STM's correctness reduces to two properties:
\begin{enumerate}
  \item Every execution is \emph{conflict-equivalent} to the sequential
    execution of the same block in the block's canonical order.
  \item Every committed write is visible to every strictly-later transaction
    that reads the same storage location.
\end{enumerate}

These properties are purely local to the block's set of transactions; they
make no reference to the fact that blocks are arranged in a linear chain.
This is the opening we exploit.

\subsection{Construction}

Let $T$ be a sequence of pending EVM transactions. The DAG-EVM builder
performs:

\begin{algorithm}[h]
\caption{BuildVertex (DAG-EVM)}
\begin{algorithmic}[1]
\State Drain up to $N$ txs from the mempool into $T$
\State Run Block-STM speculative exec on $T$ over the current state snapshot
\State Emit $(R, W) \gets$ union of per-tx read/write sets
\State $P \gets$ minimal antichain of frontier vertices whose $W$ covers $R$
\State Let $v.\mathsf{parents} \gets P$
\State $v.\mathsf{id} \gets \mathsf{keccak256}(\mathsf{txHashes} \parallel R \parallel W)$
\State Emit $v$ to the DAG engine
\end{algorithmic}
\end{algorithm}

The DAG engine (nebula) votes on acceptance. Once an antichain at rank $r$
is fully populated with accepted vertices, Quasar stamps it with a
post-quantum triple-signature certificate and $r$ becomes final. The EVM
state is updated by applying the vertices of each finalised antichain in
any valid topological order.

\subsection{Correctness Sketch}

\begin{theorem}[Non-Conflict Commutativity, Lean: \texttt{nonConflict\_commute}]
\label{thm:commute}
For any two EVM vertices $v_1, v_2$ with $\mathsf{conflicts}(v_1, v_2)
= \mathsf{false}$ and any state $\sigma$:
\[
v_2.\mathsf{apply}(v_1.\mathsf{apply}(\sigma)) = v_1.\mathsf{apply}(v_2.\mathsf{apply}(\sigma)).
\]
\end{theorem}

\begin{proof}[Proof sketch]
By case analysis on membership of each storage key $k$ in the write sets
of $v_1$ and $v_2$. The non-conflict hypothesis excludes write-write
overlap; the footprint lemmas (\texttt{writesOnly}, \texttt{readsOnly}) in
the Lean development handle the remaining four cases. Full proof in
\texttt{Consensus/DAGEVM.lean:60--135}.
\end{proof}

\begin{theorem}[Topological Order Invariance, Lean: \texttt{topo\_equivalence}]
\label{thm:topo}
Given the Siblings-Commute invariant, any two topological orderings $L_1,
L_2$ of the same accepted DAG $D$ satisfy
$\mathsf{applyList}(L_1, \sigma) = \mathsf{applyList}(L_2, \sigma)$
for every initial state $\sigma$.
\end{theorem}

\begin{theorem}[No Double Spend, Lean: \texttt{no\_double\_spend}]
If $v, w$ are both accepted, $v \neq w$, and both write the same storage
key $k$, then $v$ and $w$ must be in distinct ranks (one is an ancestor of
the other in the DAG).
\end{theorem}

\begin{corollary}[Serialisability]
Every DAG-EVM execution is conflict-equivalent to a sequential EVM
execution over the same set of transactions. An external observer sees a
linear history of state transitions; the underlying parallel vertex
topology is invisible at the state projection.
\end{corollary}

%% ============================================================================
